\documentclass[10pt,a4paper]{amsart}
\usepackage[margin=19mm]{geometry}
\usepackage{amsmath,amssymb,mathtools}
\usepackage[scr=rsfs,cal=euler]{mathalfa}
\usepackage{xfrac}
\usepackage[unicode,hidelinks,bookmarks=false]{hyperref}
\newcommand{\ie}{i.e.\ }
\newcommand{\cf}{cf.\ }
\newcommand{\resp}{respectively\ }
\newcommand{\Subsection}{\subsection}
\providecommand{\nobreakdash}{\nobreak\hbox{-}}
\setlength{\emergencystretch}{2em}
\multlinegap0pt
\usepackage{xcolor}
\usepackage{amssymb}
\newtheorem{lem}{LEMMA}
\title{Rapid stabilization of general linear systems with $F$-equivalence}
\author{Amaury Hayat, Epiphane Loko}
\date{Source: arXiv:2603.04031v1 (CC BY 4.0)}
\begin{document}
\maketitle

\noindent\emph{Source and licence.} Rapid stabilization of general linear systems with $F$-equivalence. Version arXiv:2603.04031v1 (CC BY 4.0), distributed by arXiv under the Creative Commons Attribution 4.0 International licence, http://creativecommons.org/licenses/by/4.0/, as recorded in the arXiv licence metadata for this exact version and shown on the abstract page https://arxiv.org/abs/2603.04031v1. Full licence notice: \texttt{licenses/arxiv-2603.04031-CC-BY-4.0.txt}.\par\smallskip

\noindent\textit{Editorial scope.} Counted unit: the article's lem Lemma-K\_n (source0.tex lines 1045--1047). The article's complete original proof of it is delivered with it (source0.tex lines 1052--1203, one \texttt{proof} environment of 152 lines, which the article places away from the statement). No mathematics is written, repaired or re-derived: the statement and the proof are the article's own lines, in the article's own notation, and no retained line is substituted or re-flowed.  Retained uncounted as the source-local context the proof needs: lines 718--721; lines 804--808; lines 822--828; lines 1021--1023; lines 1031--1033; lines 1039--1042.  Nothing was omitted from any retained span, so the package carries no omission record.  No retained line contains a citation, so the wrapper carries no bibliography and no bibliography entry.  No retained line contains a figure environment, an image-inclusion command or drawing code, so no image is delivered and the package declares no asset dependency.  Editorial formatting only, in addition: an AMS \texttt{amsart} preamble that restates the article's own declarations and macros used by the retained lines, namely the environments \texttt{lem} on separate counters, together with the standard packages those lines need.  The retained environments print in this document as Lemma 1 in order of appearance, whereas the article prints the same items with the numbering of its own full document.  The complete native source of the article as distributed in the arXiv e-print for this exact version is bundled unchanged as \texttt{sources/195713/source0.tex}.
\begin{equation}
\label{eq:tbb}
    T_iB_{i}= B_{i},
\end{equation}

\begin{lem}
\label{lem:lambda}
    For any $\lambda\in\mathbb{R}_{+}\setminus\mathcal{N}$, there exists $C(\lambda)>0$ such that 
    $|\lambda_n-\lambda_p+\lambda|\geq C(\lambda)|\lambda_n-\lambda_p|$.
\end{lem}

\begin{lem}
\label{lem:tech}
For any $s<\alpha-1$ we have
\begin{equation} \sum\limits_{n\in\mathbb{N}^{*}\setminus\{p\}}\frac{n^{s}}{|\lambda_{n}-\lambda_{p}+\lambda|}\lesssim p^{1- \alpha+ s}\log(p)+p^{-\alpha}, \; \forall p\in \mathbb{N}^*,
\end{equation}
where $\lesssim$ means lower or equal up to a multiplicative constant that does not depend on $p$ or $n$.
\end{lem}

\begin{lem}\label{lem:regk}
 There exists a unique sequence $(K_n)_{n\in{\mathbb{N}^*}}$ such that for any $\varepsilon \in (0,\alpha-1)$ the condition \eqref{eq:tbb} holds in $\tau(\mathcal{H}^{-\frac{1}{2}-\varepsilon})$ and $(-K_n b_n n^{-(\frac{1}{2}+\varepsilon)})_{n\in\mathbb{N}^*}\in l^2.$   
\end{lem}

    \begin{align}
   \sum_{n\in \mathbb{N}^*}-K_n b_n \tau q_n=\sum_{n\in \mathbb{N}^*}b_n \varphi_n,\label{TB=B-expression}
\end{align}  and we yet know that its right side belongs to $\tau(\mathcal{H}^{-\frac{1}{2}-\varepsilon})$ for all $\varepsilon>0$. This means, in particular, that the left-hand side does too. As a consequence $(-K_n b_n n^{-(\frac{1}{2}+\varepsilon)})_{n\in \mathbb{N}^*}\in l^2$ for all $\varepsilon \in (0,\alpha-1)$.

 \begin{equation}
 \label{K_n-expression}
    k_n:=-(K_nb_n+\lambda)\quad \forall n\in \mathbb{N}^*.
\end{equation}

  \begin{lem}\label{Lemma-K_n}
Consider $r\in (1/2-\alpha,\alpha-1/2).$ There exists $\bar\varepsilon>0$ such that the sequence $(k_n n^{\varepsilon})_{n\in \mathbb{N}^*}$ is uniformly bounded for any $\varepsilon\in (0,\bar \varepsilon).$ In particular, the sequence $(K_n b_n)_{n\in \mathbb{N}^*}$ defined by \eqref{K_n-expression} is uniformly bounded.
  \end{lem}

\begin{proof}[Proof of Lemma \ref{Lemma-K_n}]
For any $\varepsilon >0,$ it holds that $(n^{-\frac{1}{2}-\varepsilon}\lambda)_{n\in\mathbb{N}^{*}}\in l^2,$ thus using Lemma \ref{lem:regk} and \eqref{K_n-expression}, we have for all $\varepsilon\in (0,\alpha-1)$ that
$(k_n n^{-\frac{1}{2}-\varepsilon})_{n\in \mathbb{N}^*}\in l^2,$ and in particular
\begin{align} 
    (k_n n^{-r})_{n\in \mathbb{N}^*}\in l^2,\quad \forall r\in \left(\frac{1}{2},\alpha-\frac{1}{2}\right).\label{kn-regular0}
\end{align}
Still using Lemma \ref{lem:regk}, the equation \eqref{TB=B-expression} makes sense in $\tau(\mathcal{H}^{-\frac{1}{2}-\varepsilon}),$ and by expressing $K_n b_n$ and $\tau q_n$ it becomes
\begin{align*}
   \sum_{n\in \mathbb{N}^*}\lambda \sum_{p\in \mathbb{N}^*}\frac{b_{p}\varphi_p}{\lambda_n-\lambda_p+\lambda}+ \sum_{n\in \mathbb{N}^*}k_{n}\tau q_{n}= \sum_{n\in \mathbb{N}^*}b_n \varphi_n.
\end{align*}
This yields to the following
\begin{align}
   \sum_{n\in \mathbb{N}^*}\lambda \sum_{p\in \mathbb{N}^*\setminus\{n\}}\frac{b_{p}\varphi_p}{\lambda_n-\lambda_p+\lambda} = -\sum_{n\in \mathbb{N}^*}k_{n}\tau q_{n}\label{sum-regular}
\end{align}

which holds a priori in $\tau(\mathcal{H}^{-\frac{1}{2}-\varepsilon})$ for any $\varepsilon \in (0,\alpha-1).$ However, the equality \eqref{sum-regular} \textcolor{black}{now holds in a more regular space} than $\tau(\mathcal{H}^{-\frac{1}{2}-\varepsilon}).$ Indeed, if $\alpha>3/2,$ we can give a sense to \eqref{sum-regular} in $\tau(\mathcal{H}^{\varepsilon})$ for $0\leq \varepsilon\leq \alpha-3/2.$ To prove that, let us first notice by Fubini's Theorem in $\tau(\mathcal{H}^{-\frac{1}{2}-\varepsilon})$ that
\begin{align*}
   \sum_{n\in \mathbb{N}^*}\lambda \sum_{p\in \mathbb{N}^*\setminus\{n\}}\frac{b_{p}\varphi_p}{\lambda_n-\lambda_p+\lambda} = \sum_{p\in \mathbb{N}^*}b_{p}\varphi_p\lambda \sum_{n\in \mathbb{N}^*\setminus\{p\}}\frac{1}{\lambda_n-\lambda_p+\lambda}.
\end{align*}
And it holds that 
\begin{align*}
    \sum_{p\in \mathbb{N}^*}b_{p}\varphi_p\lambda \sum_{n\in \mathbb{N}^*\setminus\{p\}}\frac{1}{\lambda_n-\lambda_p+\lambda}=\sum_{p\in \mathbb{N}^*}\tau(p^{-\varepsilon}\varphi_p) \left(p^{\varepsilon}\lambda\sum_{n\in \mathbb{N}^*\setminus\{p\}}\frac{1}{\lambda_n-\lambda_p+\lambda}\right).
\end{align*}
Since $\tau(p^{-\varepsilon}\varphi_p)_{p\in \mathbb{N}^*}$ is a Riesz basis of $\tau(\mathcal{H}^{\varepsilon}),$ it suffices to show that \[\left(p^{\varepsilon}\lambda\displaystyle\sum_{n\in \mathbb{N}^*\setminus\{p\}}\frac{1}{\lambda_n-\lambda_p+\lambda}\right)_{p\in \mathbb{N}^*}\in l^2\] to conclude that the left hand side of \eqref{sum-regular} belongs to  $\tau(\mathcal{H}^{\varepsilon}).$ Using Lemma \ref{lem:lambda} and Lemma \ref{lem:tech}, we have that
\begin{align*}
    \left\|\left(p^{\varepsilon}\lambda\displaystyle\sum_{n\in \mathbb{N}^*\setminus\{p\}}\frac{1}{\lambda_n-\lambda_p+\lambda}\right)_{p\in \mathbb{N}^*}\right\|^2_{l^2}&=\sum_{p\in \mathbb{N}^*}p^{2\varepsilon}\lambda^2\left|\displaystyle\sum_{n\in \mathbb{N}^*\setminus\{p\}}\frac{1}{\lambda_n-\lambda_p+\lambda}\right|^2\\
    &\leq \frac{\lambda^2}{C(\lambda)^2}\sum_{p\in \mathbb{N}^*}p^{2\varepsilon}\left(\displaystyle\sum_{n\in \mathbb{N}^*\setminus\{p\}}\frac{1}{|\lambda_n-\lambda_p|}\right)^2\\
    &\lesssim\sum_{p\in \mathbb{N}^*} p^{2\varepsilon+2(1-\alpha)}\log^2(p),
\end{align*} 
and this converges for $\varepsilon \in (0,\alpha-3/2).$ So if $\alpha>3/2,$ the left hand side of \eqref{sum-regular} holds in $\tau(\mathcal{H}^{\varepsilon})$ for $\varepsilon \in (0,\alpha-3/2).$\\ 

Consider now the right hand side of \eqref{sum-regular}. We have that
\begin{align}
\label{eq:decompkn1}
   \sum_{n\in \mathbb{N}^*}k_{n}\tau q_{n}=\sum_{n\in \mathbb{N}^*}k_{n}n^{\varepsilon}\tau(n^{-\varepsilon} q_{n}). 
\end{align} 
Recalling that $(n^{-\varepsilon} q_{n})_{n\in \mathbb{N}^*}$ is a Riesz basis of $\mathcal{H}^{\varepsilon}$ for $\varepsilon \in (0,\alpha-3/2),$ it holds by the isomorphism property of $\tau$ that $\tau(n^{-\varepsilon} q_{n})_{n\in \mathbb{N}^*}$ is a Riesz basis of $\tau(\mathcal{H}^{\varepsilon}).$ Thus in view of \eqref{sum-regular} and \eqref{eq:decompkn1}, 
for any $\varepsilon \in (0,\alpha-3/2),$
\begin{align*}
    (k_n n^{\varepsilon})_{n\in \mathbb{N}^*}\in l^2.
\end{align*} 
In particular, we conclude that if $\alpha>3/2$, $(k_n)_{n\in \mathbb{N}^*}\in l^{\infty}.$
In view of \eqref{K_n-expression}, we have $K_n b_n=-(\lambda+k_n)$ and since $(k_n)_{n\in \mathbb{N}^*}$ belongs to $l^{\infty},$ we get that 
\begin{align}
    (K_n b_n)_{n\in \mathbb{N}^*}\in l^{\infty}\quad  \forall \alpha>3/2.\label{Kn-l-infty1}
\end{align} 
Consider now $1<\alpha\leq 3/2.$ In this case the gain of regularity in \eqref{sum-regular} is not enough to deduce the boundedness of $K_n b_n.$ However we are able to show that the equality \eqref{sum-regular} holds in $\tau(\mathcal{H}^{-\varepsilon}),$ for $\varepsilon>3/2-\alpha.$ Indeed, similarly to the previous case, the following holds
\begin{align*}
    \left\|\left(p^{-\varepsilon}\lambda\displaystyle\sum_{n\in \mathbb{N}^*\setminus\{p\}}\frac{1}{\lambda_n-\lambda_p+\lambda}\right)_{n\in \mathbb{N}^*}\right\|^2_{l^2}&=\sum_{p\in \mathbb{N}^*}p^{-2\varepsilon}\lambda^2\left|\displaystyle\sum_{n\in \mathbb{N}^*\setminus\{p\}}\frac{1}{\lambda_n-\lambda_p+\lambda}\right|^2\\
    &\leq \frac{\lambda^2}{C(\lambda)^2}\sum_{p\in \mathbb{N}^*}p^{-2\varepsilon}\left(\displaystyle\sum_{n\in \mathbb{N}^*\setminus\{p\}}\frac{1}{|\lambda_n-\lambda_p|}\right)^2\\
    &\lesssim\sum_{p\in \mathbb{N}^*} p^{-2\varepsilon+2(1-\alpha)}\log^2(p).
\end{align*} This converges for $\varepsilon>3/2-\alpha$ and then the left hand side of \eqref{sum-regular} belongs to $\tau(\mathcal{H}^{-\varepsilon}).$ 
Similarly as before and in view of \eqref{sum-regular}, 
\begin{align}
 (k_n n^{-\varepsilon})_{n\in\mathbb{N}^*}\in l^2 \quad \forall \varepsilon\in \left(\frac{3}{2}-\alpha, \alpha-\frac{1}{2}\right).\label{kn-regular1}  
\end{align}
We can see that we obtain $(\alpha-1)$ gain of regularity between \eqref{kn-regular1} and \eqref{kn-regular0}. So in the following, we will make an iterative principle to gain at each order this $(\alpha-1)$ regularity in order to prove Lemma \ref{Lemma-K_n}. So let us go back to the equation \eqref{TB=B-expression} and expressing $\tau q_n$
\begin{align*}
    \sum_{n\in \mathbb{N}^*}-K_n b_n\left(\frac{b_n\varphi_n}{\lambda}+\sum_{p\in \mathbb{N}^*\setminus\{n\}}\frac{b_p\varphi_p}{\lambda_n-\lambda_p+\lambda}\right)=\sum_{n\in \mathbb{N}^*}b_n\varphi_n.
\end{align*}
Replacing $-K_n b_n$ by $(\lambda+k_n)$ in the first term only and using Fubini's Theorem in $\tau(\mathcal{H}^{-\frac{1}{2}-\varepsilon})$ we get:
\begin{align*}
  \sum_{n\in \mathbb{N}^*}\frac{k_n b_n}{\lambda}\varphi_n-\sum_{p\in \mathbb{N}^*}b_p\varphi_p\sum_{n\in \mathbb{N}^*\setminus\{p\}}\frac{K_n b_n}{\lambda_n-\lambda_p+\lambda}=0.  
\end{align*}
Since $(\varphi_n)_{n\in \mathbb{N}^*}$ is a Riesz basis, we have by identification that
\begin{align*}
  \frac{k_m b_m}{\lambda}=b_m\sum_{n\in \mathbb{N}^*\setminus\{m\}}\frac{K_n b_n}{\lambda_n-\lambda_m+\lambda}\quad \forall m\in \mathbb{N}^*.  
\end{align*}
Thanks to the fact that $b_m\neq 0$ for all $m\in \mathbb{N}^*$ and $K_n b_n=-(\lambda+k_n),$ this implies that 
\begin{align*}
  k_m=-\lambda\sum_{n\in \mathbb{N}^*\setminus\{m\}}\frac{\lambda+k_n}{\lambda_n-\lambda_m+\lambda}\quad \forall m\in \mathbb{N}^*.  
\end{align*}
Let us now set $\lambda=e_n^0$ and $k_n=k_n^0.$ It holds that
\begin{align}
  k_m=-\lambda\sum_{n\in \mathbb{N}^*\setminus\{m\}}\frac{e_n^0+k_n^0}{\lambda_n-\lambda_m+\lambda}\quad \forall m\in \mathbb{N}^*.  \label{kn-def}
\end{align}
Thus $k_m$ can be rewritten as $k_m=e_m^1+k_m^1$ where
\begin{align*}
    e_m^1=-\lambda\sum_{n\in \mathbb{N}^*\setminus\{m\}}\frac{e_n^0}{\lambda_n-\lambda_m+\lambda}, \quad k_m^1=-\lambda\sum_{n\in \mathbb{N}^*\setminus\{m\}}\frac{k_n^0}{\lambda_n-\lambda_m+\lambda} \quad \forall m\in \mathbb{N}^*.
\end{align*}
And using Lemma \ref{lem:tech}, we have that
\begin{align}
    |e_m^1|\lesssim m^{1-\alpha}\log m+m^{-\alpha}\lesssim 1.\label{e-1}
\end{align} So, we focus on the regularity of $k_n^1.$ We define $\varepsilon_0:=\frac{3}{2}-\alpha$ and we have
\begin{align*}
    \|(m^{-\varepsilon} k_m^1)_m\|_{l^2}\lesssim \sum_{m\in \mathbb{N}^*}m^{-2\varepsilon}\left(\sum_{n\in \mathbb{N}^*\setminus\{m\}}\frac{k_n^0}{\lambda_n-\lambda_m+\lambda}\right)^2.
\end{align*}
Using Cauchy-Schwartz inequality and Lemma \ref{lem:lambda}, we get
\begin{align*}
    \|(m^{-\varepsilon} k_m^1)_m\|_{l^2}\lesssim \sum_{m\in \mathbb{N}^*}m^{-2\varepsilon}\left(\sum_{n\in \mathbb{N}^*\setminus\{m\}}\frac{|k_n^0|^2}{|\lambda_n-\lambda_m|}\right)\left(\sum_{n\in \mathbb{N}^*\setminus\{m\}}\frac{1}{|\lambda_n-\lambda_m|}\right).
\end{align*}
Then applying Lemma \eqref{lem:tech} and Fubini's \textcolor{black}{theorem}, we have
\begin{align*}
    \|(m^{-\varepsilon} k_m^1)_m\|_{l^2}&\lesssim \sum_{m\in \mathbb{N}^*}m^{-2\varepsilon+1-\alpha}\log m\left(\sum_{n\in \mathbb{N}^*\setminus\{m\}}\frac{|k_n^0|^2}{|\lambda_n-\lambda_m|}\right)\\
    &\lesssim \sum_{n\in \mathbb{N}^*}|k_n^0|^2\left(\sum_{m\in \mathbb{N}^*\setminus\{n\}}\frac{m^{-2\varepsilon+1-\alpha}\log m}{|\lambda_n-\lambda_m|}\right)\\
    &\lesssim \sum_{n\in \mathbb{N}^*}|k_n^0|^2n^{-2(\varepsilon-1+\alpha)+\sigma},
\end{align*} 
for any $\sigma>0$. As $k_n^0=k_n$ and in view of \eqref{kn-regular1}, for any $\varepsilon-(1-\alpha)\in \left(3/2-\alpha,\alpha-1/2\right)$ (i.e. $\varepsilon\in \left(5/2-2\alpha,1/2\right)$) there exists $\sigma>0$ such that this converges. Thus for $\alpha\in (1,3/2]$, the following holds
\begin{align}
    (n^{-\varepsilon} k_n^1)_{n\in \mathbb{N}^*}\in l^2, \quad \forall \varepsilon\in \left(\frac{5}{2}-2\alpha,\frac{1}{2}\right).\label{kn1-l2}
\end{align}
If $\alpha>\frac{5}{4},$ we have $\varepsilon_1=2\alpha-\frac{5}{2}>0$ and then $(n^{\delta} k_n^1)_{n\in \mathbb{N}^*}\in l^2$ for all $\delta\in[0,\varepsilon_1).$ This ensures that $(k_n^1)_{n\in \mathbb{N}^*}$ is uniformly bounded. And combining with \eqref{e-1}, it ensures that $(k_n)_{n\in \mathbb{N}^*}$ also is uniformly bounded. Therefore $(K_n b_n)_{n\in \mathbb{N}^*}$ is uniformly bounded because $K_n b_n=-(\lambda+k_n).$ \\
If $1<\alpha\leq\frac{5}{4},$ the uniform boundedness of $K_n b_n$ cannot be immediately deduced, but we need to iterate again. So using the definition \eqref{kn-def}, we have from $k_n^0=k_n,$ that
\begin{align*}
     k_m&=-\lambda\sum_{n\in \mathbb{N}^*\setminus\{m\}}\frac{e_n^0+k_n^0}{\lambda_n-\lambda_m+\lambda}\\
     &=-\lambda\sum_{n\in \mathbb{N}^*\setminus\{m\}}\frac{e_n^0}{\lambda_n-\lambda_m+\lambda}-\lambda\sum_{n\in \mathbb{N}^*\setminus\{m\}}\frac{k_n}{\lambda_n-\lambda_m+\lambda}\\
     &=-\lambda\sum_{n\in \mathbb{N}^*\setminus\{m\}}\frac{e_n^0}{\lambda_n-\lambda_m+\lambda}-\lambda\sum_{n\in \mathbb{N}^*\setminus\{m\}}\frac{1}{\lambda_n-\lambda_m+\lambda}\left(-\lambda\sum_{p\in \mathbb{N}^*\setminus\{n\}}\frac{e_p^0+k_p^0}{\lambda_p-\lambda_n+\lambda}\right).\\
\end{align*}
Referring to the definitions of $e_n^1$ and $k_n^1,$ it holds that
\begin{align}
 k_m&=e_m^1-\lambda\sum_{n\in \mathbb{N}^*\setminus\{m\}}\frac{e_n^1}{\lambda_n-\lambda_m+\lambda} -\lambda\sum_{n\in \mathbb{N}^*\setminus\{m\}}\frac{k_n^1}{\lambda_n-\lambda_m+\lambda}\nonumber\\
 &=e_m^1+e_m^2+k_m^2,\label{e1e2k2}
\end{align} where
\begin{align*}
  e_m^2=  -\lambda\sum_{n\in \mathbb{N}^*\setminus\{m\}}\frac{e_n^1}{\lambda_n-\lambda_m+\lambda}, \quad k_m^2=-\lambda\sum_{n\in \mathbb{N}^*\setminus\{m\}}\frac{k_n^1}{\lambda_n-\lambda_m+\lambda} \quad \forall m\in \mathbb{N}^*.
\end{align*}
Since $k_m=e_m^1+k_m^1,$ we have in view of \eqref{e1e2k2} that $k_m^1=e_m^2+k_m^2.$ This means that if $(e_m^2)_{m\in \mathbb{N}}$ and $(k_m^2)_{m\in \mathbb{N}}$ are uniformly bounded, then $(k_m^1)_{m\in \mathbb{N}}$ is and consequently $(k_m)_{m\in \mathbb{N}}.$ Let us first notice that combining \eqref{e-1} and Lemma \ref{lem:tech}, we have
\begin{align*}
    |e_m^2|\lesssim m^{1-\alpha}\log n+m^{-\alpha}\lesssim 1.
\end{align*}
For the $(k_m)_{m\in \mathbb{N}},$ we use again Lemma \eqref{lem:tech}  and Fubini's Theorem to have, similarly as previously:
\begin{align*}
    \|(m^{-\varepsilon} k_m^2)_m\|_{l^2}&\lesssim \sum_{m\in \mathbb{N}^*}m^{-2\varepsilon}\left(\sum_{n\in \mathbb{N}^*\setminus\{m\}}\frac{k_n^1}{\lambda_n-\lambda_m+\lambda}\right)^2\\
    &\lesssim \sum_{n\in \mathbb{N}^*}|k_n^1|^2n^{-2(\varepsilon-1+\alpha)+\sigma},
\end{align*}
for any $\sigma>0$.
In view of \eqref{kn1-l2}, there exists $\sigma>0$ such that this converges for any $\varepsilon -(1-\alpha)\in \left(\frac{5}{2}-2\alpha,\frac{1}{2}\right)$ meaning for $\varepsilon\in \left(\frac{7}{2}-3\alpha,\frac{3}{2}-\alpha\right).$ Thus, for $\alpha\in (1,3/2],$
\begin{align*}
    (n^{-\varepsilon}k_n^2)_{n\in \mathbb{N}^*}\in l^2, \quad \forall \varepsilon\in \left(\frac{7}{2}-3\alpha,\frac{3}{2}-\alpha\right).
\end{align*}

If $\alpha > 7/6,$ then $\varepsilon_2=3\alpha-\frac{7}{2}>0$ and $(n^{\delta}k_n^2)_{n\in \mathbb{N}^*}\in l^2,$ for any $\delta\in [0,\varepsilon_2).$ This ensures that $(k_n^2)_{n\in \mathbb{N}^*}$ is uniformly bounded and then $(k_n^2)_{n\in \mathbb{N}^*}$ is. Thus, $\left(K_n b_n=-(e_n^0+e^1_n+e_n^2+k_n^2)\right)_{n\in \mathbb{N}^*}$ is uniformly bounded.\\

If $1<\alpha\leq \frac{7}{6},$ we can continue the induction and have
$k_n^i=e_n^{i+1}+k_n^{i+1}$ for all $i\in \mathbb{N}$ where
\begin{align*}
  e_n^{i+1}=  -\lambda\sum_{m\in \mathbb{N}^*\setminus\{n\}}\frac{e_m^i}{\lambda_m-\lambda_n+\lambda}, \quad k_n^{i+1}=-\lambda\sum_{m\in \mathbb{N}^*\setminus\{n\}}\frac{k_m^i}{\lambda_m-\lambda_n+\lambda} \quad \forall n\in \mathbb{N}^*, i\in \mathbb{N}.
\end{align*}
And for any $n\in \mathbb{N}^*,$
\begin{align*}
    K_n b_n=-\left(k_n^i+\sum_{j=0}^{i}e_n^j\right), \quad \forall i\in \mathbb{N}
\end{align*}
Thus, based on the previous computations, we obtain by induction the following:
\begin{align}
    |e_n^i|\lesssim n^{1-\alpha}\log n+n^{-\alpha}\lesssim 1,\label{en-i}\\
    \|(n^{-\varepsilon}k_n^{i+1})_{n\in \mathbb{N}^*}\|_{l^2}\lesssim \sum_{n\in \mathbb{N}^*}|k_n^i|^2n^{-2(\varepsilon-1+\alpha)+\sigma}\log^2 n\label{nk-i+1}
\end{align} for any $\sigma>0$ and \eqref{nk-i+1} converges for $\varepsilon\in (-\varepsilon_i,-\varepsilon_{i-2})$ with $\varepsilon_i=(\alpha-1)i-\varepsilon_0$, \textcolor{black}{with $\varepsilon_{0}=1/2$,} for all $i\geq 2.$ It is worth stressing that there exists a finite $i_0\in \mathbb{N}$ such that $\varepsilon_{i_0}>0.$ In this case we have $(n^{\delta}k_n^{i_0+1})_{n\in \mathbb{N}^*}\in l^2$ for any $\delta \in [0,\varepsilon_{i_0})$ meaning that $(k_n^{i_0+1})_{n\in \mathbb{N}^*}\in l^{\infty}.$ Combining this with \eqref{en-i} we conclude that 
\begin{align}
    \left(K_n b_n=-\left(k_n^{i_0}+\sum_{j=0}^{i_0}e_n^j\right)\right)_{n\in \mathbb{N}^*}\in l^{\infty}\quad \forall \alpha\in (1, 3/2]. \label{Kn-l-infty2}
\end{align}
Therefore we conclude Lemma \ref{Lemma-K_n} from \eqref{Kn-l-infty1} and \eqref{Kn-l-infty2}.
\end{proof}

\end{document}
